\documentclass[10pt, aspectratio=169]{beamer}

\usepackage[utf8]{inputenc}
\usepackage[T1]{fontenc}
\usepackage[english]{babel}
\usepackage{amsmath, amssymb}
\usepackage{bm}
\usepackage{graphicx}
\usepackage{booktabs}

\usetheme{metropolis}
\usecolortheme{seahorse}
\setbeamertemplate{footline}[frame number]

\title{The Coupling Matters}
\subtitle{A Controlled Study of Flow Matching with OT Plans}
\author{Mael Tremouille \and Giovanni Dettori}
\institute{ENSAE Paris --- M2 Optimal Transport}
\date{\today}

\begin{document}

\maketitle

% -----------------------------------------------------------------------------
\begin{frame}{Outline}
    \tableofcontents
\end{frame}

% -----------------------------------------------------------------------------
\section{Setting}

\begin{frame}{Flow Matching in one slide}
    Train $v_\theta$ such that
    \[
        \dot x(t) = v_\theta(x(t), t)
        \quad\text{transports }\nu_0 \text{ onto } \nu_1.
    \]
    With a coupling $\pi \in \Pi(\nu_0, \nu_1)$ and the linear path
    $x_t = (1-t)\,x_0 + t\,x_1$:
    \[
        \mathcal L(\theta) = \mathbb E_{(x_0, x_1) \sim \pi,\; t}\;
        \| v_\theta(x_t, t) - (x_1 - x_0) \|^2.
    \]
    Optimum: $v^\star(x, t) = \mathbb E_\pi[\,x_1 - x_0 \mid x_t = x\,]$.
\end{frame}

\begin{frame}{Why the coupling matters}
    Two sources $x_0^{(1)} \neq x_0^{(2)}$ landing on the same $x_t$
    contribute potentially conflicting velocity targets.
    \begin{itemize}
        \item independent $\pi$: large $\Var(x_1 - x_0 \mid x_t)$;
        \item OT $\pi^\star$: paths rarely cross at the same $x_t$;
        \item Sinkhorn $\pi^\varepsilon$: smooth interpolation in $\varepsilon$.
    \end{itemize}
\end{frame}

% -----------------------------------------------------------------------------
\section{Experimental design}

\begin{frame}{Identical training, only $\pi$ changes}
    \begin{itemize}
        \item moons $\to$ 8-Gaussians; $N = 1000$.
        \item 4-layer MLP, width 128, SiLU; Adam, $T = 2000$ steps.
        \item Same MLP init, same $t$-schedule, same rollout solver.
        \item Couplings: \emph{independent}, \emph{Hungarian}, \emph{Sinkhorn-sampled},
              \emph{Sinkhorn-barycentric}, \emph{perturbed-OT($\alpha$)}.
    \end{itemize}
\end{frame}

\begin{frame}{Metrics}
    \begin{itemize}
        \item \textbf{Path Energy} $\int_0^1 \|v_\theta\|^2 dt$ and NPE.
        \item \textbf{Endpoint $W_2^2$} at NFE $\in \{4, 8, 16, 32, 64\}$.
        \item \textbf{Training proxies}: loss variance, kernel estimate of
              $\mathbb E\,\Var(x_1 - x_0 \mid x_t)$.
        \item \textbf{Geometry}: $\|v_{t+\Delta t} - v_t\|^2$, angular deviation.
        \item \textbf{Mode coverage} on 8-Gaussians.
    \end{itemize}
\end{frame}

% -----------------------------------------------------------------------------
\section{Results}

\begin{frame}{Headline}
    % \includegraphics[width=\linewidth]{../report/fig_headline.pdf}
    \begin{center}
        \emph{[insert headline figure once benchmark CSV is produced]}
    \end{center}
\end{frame}

\begin{frame}{Sinkhorn $\varepsilon$ sweep}
    \begin{center}
        \emph{[Giovanni: insert figure]}
    \end{center}
\end{frame}

\begin{frame}{Perturbed-OT $\alpha$ sweep}
    \begin{center}
        \emph{[Giovanni: insert figure]}
    \end{center}
\end{frame}

\begin{frame}{Solver vs.\ NFE}
    \begin{center}
        \emph{[Mael: insert figure from nfe\_solver\_comparison]}
    \end{center}
\end{frame}

% -----------------------------------------------------------------------------
\section{Conclusion}

\begin{frame}{Take-aways}
    \begin{itemize}
        \item OT-based couplings reduce path energy and improve low-NFE quality
              \emph{at no additional training cost}.
        \item The barycentric projection further reduces target velocity
              variance --- at the cost of a slight target-geometry bias.
        \item Sinkhorn $\varepsilon$ provides a continuous knob between exact OT
              and independent sampling.
    \end{itemize}
\end{frame}

\begin{frame}{Thank you}
    \centering
    Questions?
\end{frame}

\end{document}
